% SPDX-FileCopyrightText: 2016-2026 Lari Natri <lari.natri@iki.fi>
% SPDX-License-Identifier: GPL-3.0-or-later
%
% General-purpose typesetting and control-flow utilities.
% This file is part of the 'ulsbs' package.
%
% Load: ulsbs.sty, after the main language is selected and before style defaults.
% Requires: init registers, etoolbox, xparse, xstring, TikZ and expl3.
% Later-defined fonts, colors and melody symbols are resolved when macros run.
% State: inherits @ as a letter; does not add a group or change caller catcodes.
% Allocates circle measurement/cache registers and the global previous-paragraph
% line count; installs para/after tracking. Circle cache entries are global,
% keyed by font size only. The repeat loop uses caller-scope scratch counters.
% Loads boolmap at its original position, after the utilities it accompanies.

\ProvidesFile{ulsbs-utilities.code.tex}[Shared typesetting utilities]

% Very small font size, smaller than \tiny
\newcommand{\ulNanoSize}{%
  \ifcase\@ptsize\relax
    \@setfontsize\ulNanoSize{3.5pt}{4pt}%   10pt base
  \or
    \@setfontsize\ulNanoSize{4.5pt}{5.2pt}% 11pt base
  \or
    \@setfontsize\ulNanoSize{5pt}{6pt}%     12pt base
  \else
    \@setfontsize\ulNanoSize{3.5pt}{4pt}%   fallback
  \fi
}

% Small font size, between \scriptsize and \tiny
\newcommand{\ulMicroSize}{%
  \ifcase\@ptsize\relax
    \@setfontsize\ulMicroSize{6pt}{7pt}%     10pt base
  \or
    \@setfontsize\ulMicroSize{6.5pt}{7.5pt}% 11pt base
  \or
    \@setfontsize\ulMicroSize{7pt}{8pt}%     12pt base
  \else
    \@setfontsize\ulMicroSize{6pt}{7pt}%     fallback
  \fi
}

% Typeset #1 in scratch box zero, then select the first following branch if
% its width is zero, otherwise the second. This tests width, not token emptiness.
% Example:
%   \ulsbs@ifnowidth{\somemacro}{%
%     EMPTY
%   }{%
%     NOT EMPTY
%   }
\newcommand{\ulsbs@ifnowidth}[1]{%
  \begingroup%
  \sbox\z@{#1}%
  \expandafter\endgroup%
  \ifdim\wd\z@=\z@%
    \expandafter\@firstoftwo%
  \else%
    \expandafter\@secondoftwo%
  \fi%
}

% Disables \lineskiplimit (by setting it to -20pt) for the scope of the given
% parameter, also for \ulChordLineParams. Use this to surround a section you don't
% want lineskiplimit to affect, for example to prevent chord line to jump up if
% there are tall lyrics below it, especially when there is no chord nor beat
% mark, but only a note hint circle high above. Optional parameter allows
% setting the limit.
%
% Usage:
%   \ilsl[<lineskiplimit>]{<stuff>}
%
% Example:
%   \ilsl{\[<C#>]Älä} veljeäsi |sure
%
\def\ilsl{%
  \begingroup%
  % Make # and & \active (catcode 13) to allow them within the argument; this
  % is how \[] expects them to be.
  \catcode`\#=\active\relax%
  \catcode`\&=\active\relax%
  \ulsbs@ilsl@x
}

% Helper for \ilsl: #1 is the optional limit, #2 is long-form body content.
% Both limit assignments and active accidental catcodes end with this group.
\NewDocumentCommand{\ulsbs@ilsl@x}{O{-20pt} +m}{%
  \lineskiplimit=#1\relax%
  \renewcommand{\ulChordLineLineskipLimit}{#1}%
  #2%
  \endgroup%
}


% Gets a roman uppercase number from the number given as parameter
\newcommand{\ulRomanUpperNumber}[1]{%
  \expandafter\@slowromancap\romannumeral #1@%
}

% Set a length register #1 with value #2 in the global scope.
\newlength\ulsbs@utilities@length
\gdef\ulGlobalSetLength#1#2{%
  \begingroup
    \setlength\ulsbs@utilities@length{#2} % Local assignment to a temporary register
    \global#1=\ulsbs@utilities@length     % Global assignement to #1
                        % \relax not needed because of the following \endgroup
  \endgroup             % \ulsbs@utilities@length is restored
}

% This rewrite ensures that empty pages created with \cleardoublepage, like the
% empty pages at the end of chapters, do not have headers.
\renewcommand*{\cleardoublepage}{%
  \clearpage%
  \if@twoside%
    \ifodd\c@page%
    \else%
      \hbox{}%
      \thispagestyle{empty}%
      \newpage%
      \if@twocolumn%
        \hbox{}\newpage%
      \fi%
    \fi%
  \fi%
}

% Insert an empty line when in obeylines mode, e.g. within a song verse
\newcommand{\emptyline}{\par\hspace{0pt}\par}

% Repeal the effects of \obeylines
\newcommand\disobeylines{\catcode`\^^M=5 }%

% Produces an invisible box with the vertical size (height and depth) of
% text with highest and lowest reaching characters used in the books in the
% current font style, but having zero width. Use this to make boxes same size,
% etc.
\newcommand{\ulsbs@vphantommaxtext}{\vphantom{Êg(\flt\shrp}}

% If the currently active languages name (polyglossia) is the same as the
% first argument, the second argument is expanded; otherwise the third.
% Example \ulIfLanguage{portuguese}{Círculos}{Circles}
\protect\newcommand\ulIfLanguage[3]{%
  \IfStrEq{\languagename}{#1}{% equals
    #2%
  }{%
    #3%
  }%
}

% Creates a mark that takes no space (might overlap) immediately to the left of
% current position. The mark is given as optional parameter in square brackets.
% If the mark is not given, a right-pointing arrowhead is used. Does not leak
% style.
\newcommand{\marker}[1][\symrarrowhead]{%
  \makebox[0pt][l]{{\llap{#1}}}%
}

% Creates a playout mark that takes no space (might overlap) immediately to the
% left of current position. This similar to \pomark in Lilypond, but points to
% the right instead of down.
\newcommand{\pomarkr}{%
  \marker[\color{ulPlayoutMarkColor}\symrarrowhead{ }{ }]%
}

% Circled. Use with \circled macro. This uses TikZ and is used by melody notes.

% Circled: settings. Be careful in changing these globally, as they are used for
% printing melody note hints.
\newcommand*\ulCircledWidestSymbol{g\ulsbs@mnsharpsymbol} % widest token in the notation

\newcommand*\ulCircledHorizontalPadding{.08ex} % horizontal padding for the circle with contents

\newcommand*\ulCircledCircleExtraDiameter{.05ex} % extra diameter for the circle

\newcommand*\ulCircledCircleVerticalShift{.1ex} % move ONLY the circle (up = positive)

\newcommand*\ulCircledCircleLine{.3pt} % circle line width

\newcommand*\ulCircledCircleOpacity{.7} % circle's opacity (1=opaque)

\newsavebox{\ulsbs@circledCircleBox}
\newdimen\ulsbs@circledCircleDiameter
\newdimen\ulsbs@circledCircleRadius

% Circled: compute diameter+radius for *current* font, based on
% \ulCircledWidestSymbol. Updates global dimensions; uses dimen0 and dimen2.
\newcommand*\ulsbs@circledCircleCompute{%
  \sbox{\ulsbs@circledCircleBox}{\strut \ulCircledWidestSymbol}%
  \dimen0=\wd\ulsbs@circledCircleBox
  \dimen2=\ht\ulsbs@circledCircleBox \advance\dimen2 by \dp\ulsbs@circledCircleBox
  \ifdim\dimen2>\dimen0 \dimen0=\dimen2\fi
  \advance\dimen0 by \ulCircledCircleExtraDiameter
  \global\ulsbs@circledCircleDiameter=\dimen0
  \global\ulsbs@circledCircleRadius=\dimen0 \divide\ulsbs@circledCircleRadius by 2
}

% Circled: cache circle size per font size (keyed by \f@size).
% Creates global cache macros on first use, then assigns the current dimensions.
\newcommand*\ulsbs@circledCircleGetDims{%
  \edef\ulsbs@circled@key{\f@size}%
  \@ifundefined{ulsbs@circledCircleDiameter@\ulsbs@circled@key}{%
    \ulsbs@circledCircleCompute
    \expandafter\xdef\csname ulsbs@circledCircleDiameter@\ulsbs@circled@key\endcsname{\the\ulsbs@circledCircleDiameter}%
    \expandafter\xdef\csname ulsbs@circledCircleRadius@\ulsbs@circled@key\endcsname{\the\ulsbs@circledCircleRadius}%
  }{}%
  \ulsbs@circledCircleDiameter=\csname ulsbs@circledCircleDiameter@\ulsbs@circled@key\endcsname\relax
  \ulsbs@circledCircleRadius=\csname ulsbs@circledCircleRadius@\ulsbs@circled@key\endcsname\relax
}

% Circled: \circled[<circle color>]{<content>}
% Draws a circle around the content. Used by melody notes.
% If color is omitted, the circle uses the current color (same as surrounding text).
\NewDocumentCommand{\circled}{o m}{%
  \hspace{\ulCircledHorizontalPadding}%
  \ulsbs@circledCircleGetDims
  \tikz[baseline=(T.base)]{%
    \node[inner sep=0pt, outer sep=0pt] (T) {%
      \strut\makebox[\the\ulsbs@circledCircleDiameter][c]{#2}%
    };
    % Circle only (overlay so it doesn't affect line spacing)
    \begin{scope}[overlay]
      \IfValueT{#1}{\color{#1}}% sets stroke color; does NOT affect the text
      \draw[line width=\ulCircledCircleLine, opacity=\ulCircledCircleOpacity]
        ([yshift=\ulCircledCircleVerticalShift]T.center) circle (\the\ulsbs@circledCircleRadius);
    \end{scope}
    %% The following would redraw the content text to make it appear on top of
    %% the circle if overlapping, but it slows the processing down.
    %\node[inner sep=0pt, outer sep=0pt] (T) {%
    %  \strut\makebox[\the\ulsbs@circledCircleDiameter][c]{#2}%
    %};
  }%
  \hspace{\ulCircledHorizontalPadding}%
}

% \ulRepeatTimes{<count>}{<callable>}
%
% Invoke <callable> exactly <count> times, appending one braced, 1-based index
% argument each time. #2 is not an inline template: literal ##1 is not replaced.
% <count> must be accepted by \count@=<count>; nonpositive counts do nothing.
% Uses \count@ and \@tempcnta in the current scope, as does the callable body.
% The callable must not clobber these loop counters.
%
% Example:
%   \newcount\N \N=5
%   \newcommand{\repeated}[1]{Iteration~#1\par}
%   \ulRepeatTimes{\N}{\repeated}
\newcommand{\ulRepeatTimes}[2]{% #1 = count, #2 = callable receiving {index}
  \count@=#1\relax
  \@tempcnta=1\relax
  \@whilenum\count@>0\do{%
    #2{\the\@tempcnta}%
    \advance\@tempcnta by 1\relax
    \advance\count@ by -1\relax
  }%
}

% SPDX-FileCopyrightText: 2016-2026 Lari Natri <lari.natri@iki.fi>
% SPDX-License-Identifier: GPL-3.0-or-later
%
% Boolean-map helpers.
% This file is part of the 'ulsbs' package.
%
% Requires expl3/xparse.
% Loaded by ulsbs-utilities.code.tex at its end, before repeat-bar users.
% Maps are caller-owned property lists; this module allocates no map or lifecycle
% hooks. Expl3 operations accept a property-list variable, while document wrappers
% accept its control-sequence name without a backslash and use local scratch tls.
% Set/reset operations emit nothing; tests execute only the selected branch.

\ProvidesFile{ulsbs-boolmap.code.tex}[ULSBS boolean maps]

% % BEGIN Map: key -> boolean
%
% These helpers implement an "array"/map from arbitrary keys to boolean values.
% Internally, the map is an expl3 property list (\prop_new:N ...).
%
% Each key maps to the literal token list value: "true" or "false"
%
% Property-list keys are compared as strings under expl3 property-list rules,
% not by token identity. Avoid leading/trailing spaces in keys.
%
% Like most TeX assignments, property-list modifications are scoped:
%   - \prop_put:Nnn, \prop_clear:N are LOCAL (inside TeX groups)
%   - \prop_gput:Nnn, \prop_gclear:N are GLOBAL (persist outside groups)
%
% The helpers below come in local and global variants for setting/resetting.
%
% If a key is not present in the map, it is treated as "false".
%
% Example:
%   \prop_new:N \g_seen_prop
%   \ulsbs_boolmap_gset_true:Nn  \g_seen_prop {7}
%   \ulsbs_boolmap_if_true:NnTF  \g_seen_prop {7} {YES} {NO}
%   \ulsbs_boolmap_greset:N      \g_seen_prop
%
% Example using with LaTeX2e wrappers:
%   \ExplSyntaxOn
%   \prop_new:N \g_seen_prop
%   \ExplSyntaxOff
%
%   \ulBoolSetTrue{g_seen_prop}{7}
%   \ulIfBool{g_seen_prop}{7}{YES}{NO}
%   \ulBoolReset{g_seen_prop}
%
% Expl3 :Nn helpers do not expand the supplied key. Document wrappers fully
% expand it first; use \the for a count register, not before an ordinary macro.

\ExplSyntaxOn

% \ulsbs_boolmap_set_true:Nn <prop> {<key>}
% \ulsbs_boolmap_set_false:Nn <prop> {<key>}
%
% Set the boolean value for <key> in <prop> to true/false.
%
% LOCAL assignment: the change is undone automatically at the end of the current
%                   TeX group.
\cs_new_protected:Npn \ulsbs_boolmap_set_true:Nn  #1#2 { \prop_put:Nnn  #1 {#2} {true} }

\cs_new_protected:Npn \ulsbs_boolmap_set_false:Nn #1#2 { \prop_put:Nnn  #1 {#2} {false} }

% \ulsbs_boolmap_gset_true:Nn <prop> {<key>}
% \ulsbs_boolmap_gset_false:Nn <prop> {<key>}
%
% Set the boolean value for <key> in <prop> to true/false.
%
% GLOBAL assignment: the change persists outside groups.
\cs_new_protected:Npn \ulsbs_boolmap_gset_true:Nn  #1#2 { \prop_gput:Nnn #1 {#2} {true} }

\cs_new_protected:Npn \ulsbs_boolmap_gset_false:Nn #1#2 { \prop_gput:Nnn #1 {#2} {false} }

% \ulsbs_boolmap_reset:N <prop>
%
% Remove all key/value pairs from <prop>.
%
% LOCAL reset: only affects the current group scope. The previous contents are
%              restored at the end of the group.
\cs_new_protected:Npn \ulsbs_boolmap_reset:N  #1 { \prop_clear:N  #1 }

% \ulsbs_boolmap_greset:N <prop>
%
% Remove all key/value pairs from <prop>.
%
% GLOBAL reset: the map becomes empty globally and all keys will behave as
%               "false" (the default).
\cs_new_protected:Npn \ulsbs_boolmap_greset:N #1 { \prop_gclear:N #1 }

% \ulsbs_boolmap_if_true:NnTF <prop> {<key>} {<true-code>} {<false-code>}
% \ulsbs_boolmap_if_true:NnT  <prop> {<key>} {<true-code>}
% \ulsbs_boolmap_if_true:NnF  <prop> {<key>} {<false-code>}
%
% Conditional test for a boolean stored in <prop> under <key>.
%
% Behavior:
% - If <key> exists and its stored value is exactly "true" -> TRUE branch
% - Otherwise (value is "false" OR key is missing) -> FALSE branch
% - Missing keys are treated as false
% Uses local \l_tmpa_tl scratch for lookup; no map mutation or visible output.
\prg_new_conditional:Npnn \ulsbs_boolmap_if_true:Nn #1#2 { T, F, TF }
  {
    \prop_get:NnNTF #1 {#2} \l_tmpa_tl
      {
        \str_if_eq:VnTF \l_tmpa_tl {true}
          { \prg_return_true: }
          { \prg_return_false: }
      }
      { \prg_return_false: }
  }

% LaTeX2e wrappers

% \ulBoolSetTrue{<prop-name>}{<key>}
% Set <prop-name>[<key>] = true (GLOBAL).
% Fully expand the key into local scratch before storing the literal true value.
\NewDocumentCommand{\ulBoolSetTrue}{m m}
  {
    \tl_set:Nx \l__ulsbs_boolmap_key_tl {#2}
    \prop_gput:cVn {#1} \l__ulsbs_boolmap_key_tl {true}
  }

% \ulBoolSetFalse{<prop-name>}{<key>}
% Set <prop-name>[<key>] = false (GLOBAL).
% Fully expand the key into local scratch before storing the literal false value.
\NewDocumentCommand{\ulBoolSetFalse}{m m}
  {
    \tl_set:Nx \l__ulsbs_boolmap_key_tl {#2}
    \prop_gput:cVn {#1} \l__ulsbs_boolmap_key_tl {false}
  }

% \ulBoolReset{<prop-name>}
% Clear ALL entries from the boolean map <prop-name> (GLOBAL).
\NewDocumentCommand{\ulBoolReset}{m}
  { \prop_gclear:c {#1} }

% \ulIfBool{<prop-name>}{<key>}{<true-code>}{<false-code>}
%
% Test whether <prop-name>[<key>] is true.
%
% Behavior:
%
% - If the key exists AND its value is exactly "true" -> <true-code>
% - Otherwise (missing key OR value "false")          -> <false-code>
%
% Missing keys default to FALSE.
% Expand the key into local scratch, read without mutating the map, then execute
% #3 or #4 in the caller's scope (no implicit group around branch code).
\NewDocumentCommand{\ulIfBool}{m m m m}
  {
    \tl_set:Nx \l__ulsbs_boolmap_key_tl {#2}
    \prop_get:cVNTF {#1} \l__ulsbs_boolmap_key_tl \l__ulsbs_boolmap_val_tl
      { \str_if_eq:VnTF \l__ulsbs_boolmap_val_tl {true} {#3} {#4} }
      { #4 }
  }

\ExplSyntaxOff

% % END Map: key -> boolean

